\section{An ungraded unit, integer parity, and the complete interaction}
\label{sec:FMC}

\paragraph{Intended data and provenance.}
The user specifies a globally trivial element $\tau$ with no parity;
ordinary integers carry their intrinsic even or odd parity.
The integer $1$ is not identified with $\tau$.
The prescribed values are $\tau+\tau=\tau$, $\tau^2=\tau$,
$\tau1=1$, $0\tau=0$, $\tau+1=1$, and $1+1=2$.
The global zero is shared, and may be represented by the empty set.
The calculations below retain these values. A requested additional
distributive equality is examined explicitly rather than added
silently as an axiom. The extension of the displayed mixed rules
to all integers is specified here as a constructed completion;
its further values are not attributed to the user.

The comparison with $\mathbb F_1$ uses Alain Connes and Caterina
Consani, \href{https://arxiv.org/abs/1502.05585v2}{\emph{Absolute
algebra and Segal's Gamma sets}}, original author \nolinkurl{F1_corr.tex},
formulas \texttt{functsum}, \texttt{addition}, and propositions
\texttt{mono2sss}, \texttt{monomono}, \texttt{sssalg1}.
The finite calculations below are proved directly; no novelty claim
is made for the standard free-monoid or idempotent-ring constructions.

\subsection{Full elements and a single explicit operation pair}
Write
\begin{equation}\tag{FMC1}
I=\{n_{\bar n}:n\in\mathbb Z,\ \bar n=n\bmod2\},\qquad
X=I\sqcup\{\tau\}.
\end{equation}
The notation $n_{\bar n}$ retains the value and its intrinsic parity;
it is not a second copy of the integer for each parity. The parity
map is defined on $I$, and is undefined at $\tau$. We identify the
one integer zero with the one global zero, denoted $0$.
Replacing that single representative by $\varnothing$ through a
bijection transports every operation below and introduces no new zero.
The integer operations are
\begin{equation}\tag{FMC2}
n_{\bar n}+m_{\bar m}=(n+m)_{\overline{n+m}},\qquad
n_{\bar n}m_{\bar m}=(nm)_{\overline{nm}}.
\end{equation}
Define the remaining operation values, symmetrically, by
\begin{equation}\tag{FMC3}
\begin{gathered}
\tau x=x\tau=x\quad(x\in X),\qquad
\tau+\tau=\tau,\qquad\tau+0=0+\tau=\tau,\\
\tau+n_{\bar n}=n_{\bar n}+\tau=n_{\bar n}\quad(n\ne0).
\end{gathered}
\end{equation}
These formulas define total operations on $X$ and satisfy all the
user's displayed operation values with distinct $0$, $\tau$,
$1_{\bar1}$ and $2_{\bar0}$. There is no parity assigned to $\tau$.

Multiplication is associative and commutative, with unit $\tau$
and absorbing zero. Indeed, a product of only $\tau$ factors is
$\tau$; every other finite product is the ordinary product of its
integer factors, with parity that of the resulting integer.
Inserting or removing $\tau$ does not change it. This proves
associativity for every case, including an integer zero.
The integer $1_{\bar1}$ is the multiplicative unit on the integer
subset, but it is not the unit on $X$, since
$1_{\bar1}\tau=1_{\bar1}\ne\tau$.

Addition is commutative with identity $0$. Its associativity defect
is completely described by
\begin{equation}\tag{FMC4}
(\tau+n_{\bar n})+(-n)_{\overline{-n}}=0,\qquad
\tau+(n_{\bar n}+(-n)_{\overline{-n}})=\tau\quad(n\ne0),
\end{equation}
and the reversed placement with $\tau$ last. There are no other
cases. To verify exhaustion, triples of integers use integer
associativity; triples with two or three $\tau$ values are checked
by FMC3. With exactly one $\tau$ in the middle, the two associations
agree by FMC3, whether either integer is zero or not. With $\tau$
first or last, a zero integer drops out; two nonzero integers give
the same value unless their sum is zero, which gives FMC4.
Thus this completion is an explicit algebra, but is not an
associatively additive semiring. That limitation is a proved
property of the specified operations, not a deletion of the integer
layer or its quantum-number data.

Its entire distributivity discrepancy can also be evaluated.
For an integer multiplier $k\ne0$, the ordered pair of outputs is
\begin{equation}\tag{FMC5}
\begin{aligned}
\bigl(k(\tau+\tau),\ k\tau+k\tau\bigr)&=(k,2k),\\
\bigl(k(\tau+n),\ k\tau+kn\bigr)&=(kn,k+kn)\quad(n\ne0).
\end{aligned}
\end{equation}
Every output in FMC5 carries its resulting integer parity.
The difference of right and left integer values is exactly $k$.
All other triples satisfy the distributive equation: multiplier
$\tau$ is the identity action; multiplier zero gives zero;
two integer summands use ordinary distributivity; and the remaining
mixed case has the other summand zero and gives $k=k+0$.
This proves the complete discrepancy without inferring distributivity
from notation.

\subsection{Two exact arithmetic maps, with their full defects}
There are maps
\begin{equation}\tag{FMC6}
\begin{aligned}
r_+:X\to\mathbb Z,&\qquad r_+(n_{\bar n})=n,&r_+(\tau)=0,\\
r_\times:X\to\mathbb Z,&\qquad r_\times(n_{\bar n})=n,&r_\times(\tau)=1.
\end{aligned}
\end{equation}
The first preserves addition and zero, by checking the integer,
$\tau$--zero, $\tau$--nonzero, and $\tau$--$\tau$ pairs in FMC3.
It does not preserve multiplication: $r_+(\tau n)=n$ whereas
$r_+(\tau)r_+(n)=0$ for $n\ne0$.
The second preserves multiplication, zero, and the multiplicative
unit, by the product description following FMC3. Its complete
additive defect is
\begin{equation}\tag{FMC7}
r_\times(x)+r_\times(y)-r_\times(x+y)
=\begin{cases}
1,&\text{at least one of $x,y$ is $\tau$ and neither is $0$},\\
0,&\text{otherwise}.
\end{cases}
\end{equation}
For two integers this is zero by ordinary addition. The three
remaining pair types give $1+1-1=1$, $1+n-n=1$ for $n\ne0$,
and $1+0-1=0$. This proves every case.
Neither map identifies the elements of $X$ before its stated
application. In particular $r_\times(\tau)=1$ is a value of a map,
not an equality $\tau=1_{\bar1}$ in $X$ or an assignment of parity
to $\tau$.

\subsection{The global action and the retained occurrence count}
Let $\operatorname{End}(I,+)$ use pointwise addition and composition.
For every $x\in X$, define the endomorphism
\begin{equation}\tag{FMC8}
\rho(x)(n_{\bar n})=
(r_\times(x)n)_{\overline{r_\times(x)n}}.
\end{equation}
Integer distributivity proves each map is additive. The
multiplicativity of $r_\times$ proves
$\rho(xy)=\rho(x)\circ\rho(y)$, with
$\rho(\tau)=\operatorname{id}_I$ and $\rho(0)=0$.
The restriction to $B=\{0,\tau\}$ is injective, as its two maps
have different values at $1_{\bar1}$.
The identity action preserves intrinsic parity; this does not place
$\tau$ in the domain of the parity map.
FMC7 gives the full operator discrepancy
\begin{equation}\tag{FMC9}
\rho(x)+\rho(y)-\rho(x+y)=
\begin{cases}\operatorname{id}_I,&\text{in the first case of FMC7},\\
0,&\text{otherwise}.
\end{cases}
\end{equation}
In particular its value at $1_{\bar1}$ is $1_{\bar1}$ in the first
case. The summed action and the collapsed action themselves have
values $2_{\bar0}$ and $1_{\bar1}$ for $x=y=\tau$.

Retain every occurrence by the exact pair of maps
\begin{equation}\tag{FMC10}
B\xleftarrow{q}\mathbb N\xrightarrow{c}\operatorname{End}(I,+),
\quad q(0)=0,\quad q(k)=\tau\ (k>0),\quad
c(k)(n_{\bar n})=(kn)_{\overline{kn}}.
\end{equation}
Both are unital semiring homomorphisms. For $q$, nonnegative sums
are positive exactly when at least one summand is positive, and
products are positive exactly when both factors are positive.
These are precisely Boolean addition and multiplication on $B$.
For $c$, the statements are integer distributivity and composition
of scalar multiplication. No map, even a set map, can factor $c$
through $q$, because $q(1)=q(2)$ while
$c(1)(1_{\bar1})=1_{\bar1}\ne2_{\bar0}=c(2)(1_{\bar1})$.
For every positive count the exact discrepancy is
\begin{equation}\tag{FMC11}
c(k)-\rho(q(k))=(k-1)\operatorname{id}_I.
\end{equation}
Thus the entire missing data are specified, rather than described
by an unspecified failure of a map. No order-of-evaluation convention
makes $c$ factor through $q$.

\subsection{A universal ring receiver retaining the local integer unit}
There is a ring receiver that preserves the entire integer arithmetic
and the multiplication of $X$, while recording the failure of its
mixed addition. Define
\begin{equation}\tag{FMC12}
R=\mathbb Z[U]/(U^2-U),\quad
j(\tau)=1_R,\qquad j(n_{\bar n})=nU.
\end{equation}
Every element has a unique expression $a+bU$, and its complete
product is
\begin{equation}\tag{FMC13}
(a+bU)(c+dU)=ac+(ad+bc+bd)U.
\end{equation}
Existence of that expression follows by replacing $U^m$ with $U$
for $m\ge1$. Uniqueness follows by evaluating a zero expression at
$U=0$ and $U=1$, obtaining $a=0$ and $a+b=0$.
These evaluations give the exact ring isomorphism
\begin{equation}\tag{FMC14}
R\longrightarrow\mathbb Z\times\mathbb Z,\quad
a+bU\longmapsto(a,a+b),\qquad
(r,s)\longmapsto r+(s-r)U.
\end{equation}
Formula FMC13 or direct substitution verifies both operations and
the inverse. The integer layer is the ideal $\mathbb ZU$, with
internal unit $U$, and has its original parity $nU\mapsto n\bmod2$.
The global unit $1_R$ is not in that ideal, so this parity is not
defined on it. In particular $j$ is injective on $X$: its integer
images have first coordinate zero in FMC14, while $j(\tau)$ has
first coordinate one. It preserves zero and multiplication and
all sums of integers, using $U^2=U$. Its remaining additive defects
are exactly $1_R$ in the first case of FMC7 and zero otherwise.

For any commutative ring $A$ receiving a multiplicative map from $X$
that preserves zero and the global unit and is additive on the
integer subset, the image $u$ of $1_{\bar1}$ satisfies $u^2=u$,
and the image of $n_{\bar n}$ is $nu$. Therefore the unique unital
ring map $R\to A$ carrying this data is $a+bU\mapsto a1_A+bu$.
FMC13 proves multiplicativity; addition and uniqueness follow
from the displayed expressions. This establishes the precise
universal property of this receiver. Requiring it also to preserve
$\tau+\tau=\tau$ forces $1_A+1_A=1_A$, hence $1_A=0_A$.

The ring spectrum of this receiver is fully determined rather than
inferred to be the original support spectrum. Under FMC14, its
primes are exactly
\begin{equation}\tag{FMC15}
\mathfrak p\times\mathbb Z\quad\hbox{and}\quad
\mathbb Z\times\mathfrak p\qquad
(\mathfrak p\in\operatorname{Spec}\mathbb Z).
\end{equation}
Indeed $(1,0)(0,1)=0$ forces any prime to contain one of the two
idempotents, and it cannot contain both because their sum is the
unit. In the quotient by the contained component, the prime is
exactly a prime of the other copy of $\mathbb Z$. Conversely every
listed ideal is prime by projection. The two idempotent basic
opens are disjoint, exhaust the spectrum, and each identifies
with $\operatorname{Spec}\mathbb Z$ through projection. Thus this
particular receiver has two clopen arithmetic components. It is
not silently substituted for the programme's ordinal-sum support
spectrum. The exact comparison with that spectrum is proved in FGM.

\subsection{Cohomology, tensor powers, and the scope of the unit action}
The scalar receiver in FMC8 also specifies exactly what the new
pointed unit does to every original arithmetic complex on which
it acts by scalar extension. For a complex $(C^\bullet,d)$ of
abelian groups or vector spaces its degreewise action is
\begin{equation}\tag{FMC16}
\rho(\tau)|_{C^r}=\operatorname{id}_{C^r},\qquad
d\rho(\tau)=\rho(\tau)d=d.
\end{equation}
Both composites send an element $x$ to $dx$, proving the cochain
identity. It restricts to the identity on $\ker d$ and on
$\operatorname{im}d$, so the induced map on
$H^r=\ker d/\operatorname{im}d$ is exactly the identity.
For a chain map $f:C\to D$, the usual cone has components
$D^r\oplus C^{r+1}$ and differential
$(y,x)\mapsto(d_Dy+f(x),-d_Cx)$. Applying the two identity
actions before or after this full differential gives the same
pair, including its minus sign. Thus the same assertion holds
on that cone and its cohomology. It makes no claim that a new
nonadditive source has already been identified with this complex.

For any operator $F$ on the receiving vector space and every
positive tensor order $k$,
\begin{equation}\tag{FMC17}
\rho(\tau)F=F\rho(\tau)=F,\qquad
\rho(\tau)^{\otimes k}=\operatorname{id},\qquad
(\rho(\tau)F)^{\otimes k}=F^{\otimes k}.
\end{equation}
The tensor equality follows by evaluation on every pure tensor;
pure tensors generate the algebraic tensor product, so it holds
on all elements. In particular, for an eigenvector $Fv=\alpha v$,
\begin{equation}\tag{FMC18}
F^{\otimes k}(v^{\otimes k})=\alpha^k v^{\otimes k}.
\end{equation}
Every factor is displayed: inserting this unit changes none of
the eigenvalues or their tensor powers. This locates exactly
what the unit action does in a tensor-amplification calculation.
An additional weight bound does not follow from inserting it.
Any different effect must come from the additional geometric
maps, not from a different value for this identity action.

For commuting scalar truth actions, the constructive additive
comparison instead retains the overlap:
\begin{equation}\tag{FMC19}
\rho(b\vee c)=\rho(b)+\rho(c)-\rho(b\wedge c),
\qquad b,c\in\{0,\tau\}.
\end{equation}
When both are $\tau$ it reads $I=I+I-I$; when one is zero
it reads $I=I+0-0$; when both are zero it reads $0=0$.
These cases prove the formula without assuming additivity of
$\rho$ on Boolean sums. The full lattice form is proved in FSI
and the original arithmetic trace evaluation in FTR.

\paragraph{What the construction establishes.}
FMC1--19 provide a complete non-distributive interaction, its exact
associativity and distributivity discrepancies, a global operator
action, the full occurrence-count receiver, and a universal ring
receiver with all its prime points. Every integer value and its
intrinsic parity is retained; $\tau$ never receives parity. These
calculations do not supply a sign for the original Weil pairing,
a Deligne purity estimate, or a proof or disproof of RH.
